\documentclass[12pt]{ca-capstone}
\geometry{margin=25mm}
\usepackage{setspace}
\setstretch{1.1}
\capstonerefs{knuth-taocp2,wiki-lfsr}
\begin{document}
\capstonetitle{Криптоанализ ЛКГ и РСЛОС}
\begin{center}\textbf{Теоретическая справка для преподавателя. Не выдавать студентам.}\end{center}

\section{ЛКГ и период}
ЛКГ задаётся $x_{i+1}=ax_i+c\bmod m$. Конечная последовательность состояний в конечном итоге
входит в цикл, возможно после хвоста. Таблица «состояние $\to$ первый индекс» даёт и длину
хвоста, и период; ограничение числа итераций обязательно сообщается как ошибка поиска, а не
как ответ.

Условия Халла--Добелла полного периода $m$: $\gcd(c,m)=1$; $a-1$ делится на каждый простой
делитель $m$; если $4\mid m$, то $a\equiv1\pmod4$. Для $m=2^k$ это сводится к нечётным $a,c$ и
$a\equiv1\pmod4$ при $k\ge2$.

\section{Восстановление ЛКГ}
При известных $(a,c,m)$ следующее состояние вычисляется одним шагом. Для полных выходов введём
$t_i=x_{i+1}-x_i$ и
\[
 u_i=t_{i+2}t_i-t_{i+1}^2.
\]
Истинный модуль делит каждый $u_i$, поэтому $m\mid\gcd(u_i)$. В общем случае равенство может не
выполняться, а небратимость $t_0$ создаёт неоднозначность. В работе допустимы только примеры,
для которых $m=\gcd(u_i)$ и $\gcd(t_0,m)=1$:
\[
 a=t_1t_0^{-1}\pmod m,\qquad c=x_1-ax_0\pmod m.
\]
Все переходы затем проверяются. Векторы отбираются по этим условиям, перечисление делителей и
неоднозначных троек не требуется.

\section{РСЛОС и многочлены}
Для коэффициентов $c_1,\ldots,c_L$ и состояния $s_0,\ldots,s_{L-1}$ используется
\[
 s_n=\sum_{j=1}^{L}c_j s_{n-j}\pmod2.
\]
Многочлен связей $C(x)=1+c_{1}x+\cdots+c_{L}x^{L}$, характеристический
$f(x)=x^{L}C(1/x)=x^{L}+c_{1}x^{L-1}+\cdots+c_{L}$. При вырожденном префиксе степень текущего $C$ может
быть меньше текущей длины состояния $L$: ведущие нули коэффициентов важны для восстановления
исходной последовательности.

Период и предпериод РСЛОС определяются обходом состояний $(s_n,\ldots,s_{n+L-1})$, а не только
выходных бит. При вырожденной рекуррентности возможен хвост до входа в цикл; нулевое состояние
образует цикл длины~$1$.

\section{Берлекэмп--Месси}
Хранятся текущий и предыдущий многочлены $C,B$, линейная сложность $\ell$ и индекс $m$ последнего
её увеличения. На шаге $N$ вычисляется невязка
\[
 d=s_N+\sum_{i=1}^{\ell}c_i s_{N-i}.
\]
При $d=1$ выполняется $C\gets C+x^{N-m}B$; если $2\ell\le N$, дополнительно
$\ell\gets N+1-\ell$, $B\gets C_{\mathrm{old}}$, $m\gets N$. Нулевая последовательность имеет
$\ell=0$. После $2\ell$ бит кратчайшая рекуррентность в не вырожденном случае однозначна;
работа ограничена такими входами.

Проверка состоит из сравнения степени и $f(x)$ с Sage, генерации всей последовательности по
восстановленным состоянию и коэффициентам и предсказания следующего бита. Для крайне вырожденных
префиксов соглашения о минимальном многочлене могут отличаться; benchmark использует случаи, где
курсовая трассировка и Sage согласуются.

\clearpage
\section{Поточное шифрование и контрольный список}
Шифрование $y_i=x_i\oplus k_i$ обратно само себе: $x_i=y_i\oplus k_i$. В атаке по известному
открытому тексту префикс открытого текста даёт префикс гаммы $k_i=x_i\oplus y_i$. По нему
восстанавливаются параметры ЛКГ или РСЛОС, затем генерируется продолжение гаммы и расшифровывается
оставшаяся часть сообщения. Повторное шифрование с восстановленной гаммой обязано вернуть
шифротекст. Тесты включают полный и неполный период ЛКГ, нулевой множитель, однозначное
восстановление, периоды и предпериоды РСЛОС, нулевое и вырожденное состояния, границу $2L$,
преобразование многочленов и обе атаки по известному открытому тексту.

\references
\end{document}
